\chapter*{خاتمة السلسلة}
\addcontentsline{toc}{chapter}{خاتمة السلسلة}

تنتهي هذه السلسلة عند اللقاء الأربعين، وقد دارت على سورة الفاتحة وسورة البقرة في قراءة تفسيرية تعليمية جامعة، تجمع بين تحرير مقاصد \rawi{الطبري}، وبيان مأخذ أقواله، والتنبيه على دقائق منهجه، وربط ذلك بالفقه والتربية وآداب النظر.

ومن أوضح ما تكشفه هذه الدروس أن قراءة «تفسير الطبري» ليست رحلة في جمع الأقوال فحسب، بل هي مدرسة كاملة في:
\begin{itemize}
    \item فهم لسان العرب في خدمة المعنى.
    \item جمع الأقوال قبل الترجيح بينها.
    \item تضييق دعوى النسخ، وتقديم الجمع ما أمكن.
    \item تعظيم السياق، وسبب النزول، وخاتمة الآية، والقراءة، وأثر ذلك كله في الفهم.
    \item وصل العلم بالعمل، والتفسير بالتزكية، والفقه بمراقبة القلب.
\end{itemize}

ولعل من أجل ثمرات هذه السلسلة أنها تعلم القارئ كيف يقرأ الإمام المحقق: لا يقتصر على النتيجة المجردة، ولا يتعجل في الاعتراض، ولا يذيب الفروق الدقيقة بين الأقوال، بل يصبر على طريقة العالم حتى يفهم أصوله ومآخذه.

وقد ألحقت بهذه النسخة ملاحق تجمع ما تناثر في الدروس من فوائد تفسيرية ومنهجية وفقهية وتربوية، لتكون عوناً على المراجعة السريعة، ودليلاً على أن هذه السلسلة لم تكن شرحاً تفسيرياً فحسب، بل كانت بناءً متدرجاً لملكة القراءة والتحرير والنظر.

ونسأل الله أن يبارك في الشيخ، وفي هذا الجهد، وأن يجعل هذه النسخة عوناً على فهم كتابه، ومحبة كلام أئمة الإسلام، وحسن الانتفاع به، وأن يرزق قارئها الإخلاص، والبصيرة، والعمل.
